%%--------------------------------------------------------------------
%1--------------------------------------------------------------------
\begin{glyph}{Public (公)}{public}\label{glyph:public}
Public.\\\hline
In many instances, this character serves as the opposite to 私 (p. \pageref{glyph:private}).
\end{glyph}
\gindex{Public}{公}\strindex{4}{公}

\begin{comps}
\comprtwocone&八 + ム\\\hline
\comprthreecone&Eight (Volcano) + Private\\\hline
\end{comps}

\begin{remglyph}
\hooglig{Public} \remcom{volcanoes} are being \remcom{privatized}.\\\hline
\end{remglyph}

\begin{prons}
\pronsrtwocone&\textbf{コウ (K\={O})}&---\\\hline
\pronsrthreecone&ク (KU)&おおやけ (\={o}yake)\\\hline
\end{prons}
\rindex{KOU@K\={O}}{公}
\rindex{KU}{公}
\rindex{ooyake@\={o}yake}{公}

\begin{remprons}
\hooglig{Publicly} \comon{cause} \ucomon{cooking} of \ucomkun{ore yuck}.\\\hline
\end{remprons}

%{\middel}
% &  ()\\\hline
\begin{somme}
公有&public + have = \textit{publicly owned}&コウ{\middel}ユウ \mbox{(K\={O}{\middel}Y\={U})}\\\hline
公安&public + ease = \textit{public peace}&コウ{\middel}アン \mbox{(K\={O}{\middel}AN)}\\\hline
公開&public + open = \textit{open to the public}&コウ{\middel}カイ \mbox{(K\={O}{\middel}KAI)}\\\hline
公金&public + gold = \textit{public funds}&コウ{\middel}キン \mbox{(K\={O}{\middel}KIN)}\\\hline
公道&public + road = \textit{public road}&コウ{\middel}ドウ \mbox{(K\={O}{\middel}D\={O})}\\\hline
公会&public + meet = \textit{public meeting}&コウ{\middel}カイ \mbox{(K\={O}{\middel}KAI)}\\\hline
公園&public + garden = \textit{public park}&コウ{\middel}エン \mbox{(K\={O}{\middel}EN)}\\\hline
\notyet{公平}&public + flat = \textit{fairness}&コウ{\middel}ヘイ \mbox{(K\={O}{\middel}HEI)}\\\hline
\end{somme}

\begin{decomp}{4}
ここ&は&公道&です。\\\hline
ここ&は&コウドウ&です\\\hline
this&は&public road&be/is\\\hline
\multicolumn{4}{|r|}{\textit{This is a public road.}}\\\hline
\end{decomp}
\end{multicols}
%%--------------------------------------------------------------------
%2--------------------------------------------------------------------
\begin{multicols}{2}
\begin{glyph}{Cloth (巾)}{cloth}\label{glyph:cloth}
Cloth.  Rag.\\\hline
This character frequently appears as a component in other 漢字.
\end{glyph}
\gindex{Cloth}{巾}\strindex{3}{巾}

\begin{comps}
\comprtwocone&巾\\\hline
\comprthreecone&Cloth\\\hline
\end{comps}

\begin{remglyph}
Part of the 漢字 radicals (\#{50}).\\\hline
\end{remglyph}

\begin{prons}
\pronsrtwocone&\textbf{キン (KIN)}&---\\\hline
\pronsrthreecone&はば (haba)&---\\\hline
\end{prons}
\rindex{KIN}{巾}
\rindex{haba}{巾}

\begin{remprons}
\hooglig{Clothe} the \comon{kindred} with \ucomkun{habaneros}.\\\hline
{\warn}\textit{Habaneros} is a hot variety of chilli peppers.\\\hline
\end{remprons}

%{\middel}
% &  ()\\\hline
\begin{somme}
巾&\textit{cloth, rag}&キン \mbox{(KIN)}\\\hline
手巾&hand + cloth = \textit{handkerchief}&シュ{\middel}キン \mbox{(SHU{\middel}KIN)}\\\hline
{\iscov}三角巾&triangle + cloth = \textit{sling}&サン{\middel}カク{\middel}キン \mbox{(SAN{\middel}KAKU{\middel}KIN)}\\\hline
\notyet{巾着}&cloth + wear = \textit{pouch}&キン{\middel}チャク \mbox{(KIN{\middel}CHAKU)}\\\hline
\notyet{雑巾}&miscellaneous + cloth = \textit{house-cloth}&ゾウ{\middel}キン \mbox{(Z\={O}{\middel}KIN)}\\\hline
\notyet{雑巾掛け}&dust cloth + hang = \textit{dirty work}&ゾウ{\middel}キン{\wsplit}が{\middel}け \mbox{(Z\={O}{\middel}KIN{\wsplit}ga{\middel}ke)}\\\hline
\notyet{頭巾}&head + cloth = \textit{hood}&ズ{\middel}キン \mbox{(ZU{\middel}KIN)}\\\hline
\notyet{黒頭巾}&black + hood = \textit{black hood}&くろ{\middel}ズ{\middel}キン \mbox{(kuro{\middel}ZU{\middel}KIN)}\\\hline
\end{somme}

\begin{decomp}{6}
あそこ&に&巾着&を&掛けて&下さい。\\\hline
あそこ&に&キンチャク&を&かけて&ください\\\hline
over there&に&pouch&を&hang&please\\\hline
\multicolumn{6}{|r|}{\textit{Please hang the pouch over there.}}\\\hline
\end{decomp}
\end{multicols}
\pagebreak
%%--------------------------------------------------------------------
%3--------------------------------------------------------------------
\begin{multicols}{2}
\begin{glyph}{City (市)}{city}\label{glyph:city}
City.\\\hline
A secondary meaning relates to \textit{market}.  In this sense the character can refer to anything from a \textit{simple fair} to the idea of a \textit{general business market}.\\\hline
This will be the first of four 漢字 in this chapter connected with \textit{groups of people living in a defined area}; moving from larger to smaller units will provide a sense of their relative sizes.
\end{glyph}
\gindex{City}{市}\strindex{5}{市}

\begin{comps}
\comprtwocone&亠 + 巾\\\hline
\comprthreecone&Top/Lid + Cloth\\\hline
\end{comps}

\begin{remglyph}
\hooglig{The city} on \remcom{top} has \remcom{clothes}.\\\hline
\end{remglyph}

\begin{prons}
\pronsrtwocone&\textbf{シ (SHI)}&---\\\hline
\pronsrthreecone&--- &いち (ichi)\\\hline
\end{prons}
\rindex{SHI}{市}
\rindex{ichi}{市}

\begin{remprons}
\hooglig{City} \comon{shielded} from \ucomkun{itchiness}.\\\hline
\end{remprons}

%{\middel}
% &  ()\\\hline
\begin{somme}
市内&city + inside = \textit{within the city}&シ{\middel}ナイ \mbox{(SHI{\middel}NAI)}\\\hline
市立&city + stand = \textit{municipal}&シ{\middel}リツ \mbox{(SHI{\middel}RITSU)}\\\hline
市有&city + have = \textit{city-owned}&シ{\middel}ユウ \mbox{(SHI{\middel}Y\={U})}\\\hline
\notyet{都市}&metropolis + city = \textit{town; city}&ト{\middel}シ \mbox{(TO{\middel}SHI)}\\\hline
\notyet{市民}&city + people = \textit{citizen}&シ{\middel}ミン \mbox{(SHI{\middel}MIN)}\\\hline
\notyet{市場}&city + place = \textit{marketplace}&いち{\middel}ば \mbox{(ichi{\middel}ba)}\\\hline
\notyet{都市計画}&city + plan = \textit{urban planning}&ト{\middel}シ{\middel}ケイ{\middel}カク \mbox{(TO{\middel}SHI{\middel}KEI{\middel}KAKU)}\\\hline
\end{somme}
\end{multicols}

\begin{decomp}{8}
夜&は&一人&で&市内&を&歩くないで&下さい。\\\hline
よる&は&ひとり&で&シナイ&を&あるくないで&ください\\\hline
night&は&1 person&で&within the city&を&don't walk&please\\\hline
\multicolumn{8}{|r|}{\textit{Please don't walk alone at night in the city.}}\\\hline
\end{decomp}
%%--------------------------------------------------------------------
%4--------------------------------------------------------------------
\begin{multicols}{2}
\begin{glyph}{Primary (主)}{primary}\label{glyph:primary}
Primary.\\\hline
This character conveys a sense of \textit{importance} and \textit{authority}; it is used with both concrete objects and abstract ideas.\\\hline
Note the difference between this 漢字 and 玉 (p. \pageref{glyph:jewel}).
\end{glyph}
\gindex{Primary}{主}\strindex{5}{主}

\begin{comps}
\comprtwocone&丶 + 王\\\hline
\comprthreecone&Dot + King\\\hline
\end{comps}

\begin{remglyph}
\hooglig{The primary} \remcom{jewel} of the \remcom{king}.\\\hline
\end{remglyph}

\begin{prons}
\pronsrtwocone&\textbf{シュ (SHU)}&\textbf{ぬし (nushi)}\\\hline
\pronsrthreecone&ス (SU)&おも (omo)\\\hline
\end{prons}
\rindex{SHU}{主}
\rindex{nushi}{主}
\rindex{SU}{主}
\rindex{omo}{主}

\begin{remprons}
\hooglig{Primarily} made from \comkun{new sheep}, this \comon{shooting} \ucomon{suit} can be placed in \ucomkun{omo}.\\\hline
{\warn}\textit{Omo} is a laundry detergent used for stain removal.\\\hline
\end{remprons}

%{\middel}
% &  ()\\\hline
\begin{somme}
主力&primary + strength = \textit{main force}&シュ{\middel}リョク \mbox{(SHU{\middel}RYOKU)}\\\hline
主人公&primary + person + public = \textit{main character}&シュ{\middel}ジン{\middel}コウ \mbox{(SHU{\middel}JIN{\middel}K\={O})}\\\hline
売主&sell + primary = \textit{seller}&うり{\middel}ぬし \mbox{(uri{\middel}nushi)}\\\hline
持主&hold + primary = \textit{owner}&もち{\middel}ぬし \mbox{(mochi{\middel}nushi)}\\\hline
\notyet{主食}&primary + food = \textit{staple food}&シュ{\middel}ショク \mbox{(SHU{\middel}SHOKU)}\\\hline
\notyet{主語}&primary + word = \textit{subject}&シュ{\middel}ゴ \mbox{(SHU{\middel}GO)}\\\hline
\notyet{主題}&primary + topic = \textit{theme; motif}&シュ{\middel}ダイ \mbox{(SHU{\middel}DAI)}\\\hline
\notyet{主治医}&primary + reign + doctor = \textit{attending physician}&シュ{\middel}ジ{\middel}イ \mbox{(SHU{\middel}JI{\middel}I)}\\\hline
\end{somme}

\begin{decomp}{6}
主人公&は&有名&な&美人&です。\\\hline
シュジンコウ&は&ユウメイ&な&ビジン&です\\\hline
main character&は&famous&な&beautiful woman&be/is\\\hline
\multicolumn{6}{|r|}{\textit{The main character is a famously beautiful person.}}\\\hline
\end{decomp}
\end{multicols}
\pagebreak
%%--------------------------------------------------------------------
%5--------------------------------------------------------------------
\begin{multicols}{2}
\begin{glyph}{Block (丁)}{block}\label{glyph:block}
The main sense is of a city \textit{block}.\\\hline
More obscure meanings can range from \textit{pages}, to \textit{polite}.\\\hline
This character often shows up as a component in other 漢字.
\end{glyph}
\gindex{Block}{丁}\strindex{2}{丁}

\begin{comps}
\comprtwocone&丁\\\hline
\comprthreecone&Block\\\hline
\end{comps}

\begin{remglyph}
\hooglig{Blocks} have \remcom{T-junctions}.\\\hline
\end{remglyph}

\begin{prons}
\pronsrtwocone&\textbf{チョウ (CH\={O})}&---\\\hline
\pronsrthreecone&テイ (TEI)&---\\\hline
\end{prons}
\rindex{CHOU@CH\={O}}{丁}
\rindex{TEI}{丁}

\begin{remprons}
\hooglig{Blocks} full of \ucomon{taping} \comon{chores}.\\\hline
\end{remprons}

%{\middel}
% &  ()\\\hline
\begin{somme}
丁目&block + eye = \textit{city block area}&チョウ{\middel}め \mbox{(CH\={O}{\middel}me)}\\\hline
二丁目&two + block + eye = \textit{City Block 2}&ニ{\middel}チョウ{\middel}め \mbox{(NI{\middel}CH\={O}{\middel}me)}\\\hline
四丁目&four + block + eye = \textit{City Block 4}&よん{\middel}チョウ{\middel}め \mbox{(yon{\middel}CH\={O}{\middel}me)}\\\hline
\notyet{丁重}&polite + heavy = \textit{polite; hospitable}&テイ{\middel}チョウ \mbox{(TEI{\middel}CH\={O})}\\\hline
\notyet{横丁}&sideways + block = \textit{side street}&よこ{\middel}チョウ \mbox{(yoko{\middel}CH\={O})}\\\hline
\end{somme}
\end{multicols}

\begin{decomp}{8}
私&の&新しい&家&は&四丁目&に&あります。\\\hline
わたし&の&あたらしい&いえ&は&よんチョウメ&に&あります\\\hline
I&の&new&house&は&City Block 4&に&is\\\hline
\multicolumn{8}{|r|}{\textit{My new house is in City Block 4.}}\\\hline
\end{decomp}
%%--------------------------------------------------------------------
%6--------------------------------------------------------------------
\begin{multicols}{2}
\begin{glyph}{Friend (友)}{friend}\label{glyph:friend}
Friend.
\end{glyph}
\gindex{Friend}{友}\strindex{4}{友}

\begin{comps}
\comprtwocone&ナ + 又\\\hline
\comprthreecone&Ten variant/(NA) + Again\\\hline
\end{comps}

\begin{remglyph}
\hooglig{Friends} will watch \remcom{Naruto} \remcom{again}.\\\hline
{\warn}\textit{Naruto} is an アニメ about a young ninja seeking recognition.\\\hline
\end{remglyph}

\begin{prons}
\pronsrtwocone&\textbf{ユウ (Y\={U})}&\textbf{とも (tomo)}\\\hline
\pronsrthreecone&---&---\\\hline
\end{prons}
\rindex{YUU@Y\={U}}{友}
\rindex{tomo}{友}

\begin{remprons}
\hooglig{Friend}, let's \comon{youtube} \comkun{tomorrow}.\\\hline
\end{remprons}

%{\middel}
% &  ()\\\hline
\begin{somme}
友達&friend + achieve = \textit{friend(s)}&とも{\middel}だち \mbox{(tomo{\middel}dachi)}\\\hline
友人&friend + person = \textit{friend}&ユウ{\middel}ジン \mbox{(Y\={U}{\middel}JIN)}\\\hline
友好&friend + like = \textit{friendship}&ユウ{\middel}コウ \mbox{(Y\={U}{\middel}K\={O})}\\\hline
交友&mix + friend = \textit{acquaintance}&コウ{\middel}ユウ \mbox{(K\={O}{\middel}Y\={U})}\\\hline
女友達&woman + friend = \textit{female friend}&おんなともだち \mbox{(onna{\middel}tomo{\middel}dachi)}\\\hline
親友&parent + friend = \textit{close friend}&シン{\middel}ユウ \mbox{(SHIN{\middel}Y\={U})}\\\hline
真の友&truth + friend = \textit{true friend}&シン{\middel}の{\middel}とも \mbox{(SHIN{\middel}no{\middel}tomo)}\\\hline
\notyet{悪友}&evil + friend = \textit{bad company}&アク{\middel}ユウ \mbox{(AKU{\middel}Y\={U})}\\\hline
\notyet{遊び友達}&play + friend = \textit{playmate}&あそ{\middel}び{\wsplit}とも{\middel}だち \mbox{aso{\middel}bi{\wsplit}tomo{\middel}dachi}\\\hline
\notyet{茶飲み友達}&tea drinking + friend = \textit{tea-drinking companion}&チャ{\wsplit}の{\middel}み{\wsplit}とも{\middel}だち CHA{\wsplit}no{\middel}mi{\wsplit}tomo{\middel}dachi\\\hline
\end{somme}
\end{multicols}

\begin{decomp}{7}
西山さん&の&友人&は&南米&から&来ました。\\\hline
にしやまさん&の&ユウジン&は&ナンベイ&から&きました\\\hline
Nishiyama-san&の&friend&は&South America&from&came\\\hline
\multicolumn{7}{|r|}{\textit{Nishiyama-san's friend came from South America.}}\\\hline
\end{decomp}
\pagebreak
%%--------------------------------------------------------------------
%7--------------------------------------------------------------------
\begin{multicols}{2}
\begin{glyph}{Group (団)}{group}\label{glyph:group}
Group.
\end{glyph}
\gindex{Group}{団}\strindex{6}{団}

\begin{comps}
\comprtwocone&囗 + 寸\\\hline
\comprthreecone&Enclosure/Box/Prison + Inch\\\hline
\end{comps}

\begin{remglyph}
\hooglig{The group} in the \remcom{prison} is \remcom{tiny}.\\\hline
\end{remglyph}

\begin{prons}
\pronsrtwocone&\textbf{ダン (DAN)}&---\\\hline
\pronsrthreecone&トン (TON)&---\\\hline
\end{prons}
\rindex{DAN}{団}
\rindex{TON}{団}

\begin{remprons}
\hooglig{Group} \comon{dunking} is a \ucomon{tonic}.\\\hline
\end{remprons}

%{\middel}
% &  ()\\\hline
\begin{somme}
一団&one + group = \textit{a group}&イチ{\middel}ダン \mbox{(ICHI{\middel}DAN)}\\\hline
団体&group + body = \textit{organisation}&ダン{\middel}タイ \mbox{(DAN{\middel}TAI)}\\\hline
団子&group + child = \textit{dumpling}&ダン{\middel}ご \mbox{(DAN{\middel}go)}\\\hline
入団&enter + group = \textit{joining (a group)}&ニュウ{\middel}ダン \mbox{(NY\={U}{\middel}DAN)}\\\hline
\notyet{集団}&gather + group = \textit{crowd; mass}&シュウ{\middel}ダン \mbox{(SH\={U}{\middel}DAN)}\\\hline
\notyet{団結}&group + bind = \textit{unity; solidarity}&ダン{\middel}ケツ \mbox{(DAN{\middel}KETSU)}\\\hline
\notyet{団地}&group + place = \textit{apartment complex}&ダン{\middel}チ \mbox{(DAN{\middel}CHI)}\\\hline
\notyet{記者団}&reporter + group = \textit{press organisation}&キ{\middel}シャ{\middel}ダン \mbox{(KI{\middel}SHA{\middel}DAN)}\\\hline
\notyet{代表団}&delegate + group = \textit{delegation}&ダイ{\middel}ヒョウ{\middel}ダン \mbox{(DAI{\middel}HY\={O}{\middel}DAN)}\\\hline
布団{\notcov}&spread/linen + group = \textit{bedding}&フ{\middel}トン \mbox{(FU{\middel}TON)}\\\hline
\end{somme}
\end{multicols}

\begin{decomp}{6}
内田さん&は&団子&が&好き&です。\\\hline
うちださん&は&ダンご&が&すき&です\\\hline
Uchida-san&は&dumplings&が&likes&be/is\\\hline
\multicolumn{6}{|r|}{\textit{Uchida-san likes dumplings.}}\\\hline
\end{decomp}
%%--------------------------------------------------------------------
%8--------------------------------------------------------------------
\begin{multicols}{2}
\begin{glyph}{Every (毎)}{every}\label{glyph:every}
Every.
\end{glyph}
\gindex{Every}{毎}\strindex{6}{毎}

\begin{comps}
\comprtwocone&𠂉 + 毋\\\hline
\comprthreecone&Person + Do Not/Mother\\\hline
\end{comps}

\begin{remglyph}
\hooglig{Every} \remcom{person} has a \remcom{mother}.\\\hline
\end{remglyph}

\begin{prons}
\pronsrtwocone&\textbf{マイ (MAI)}&---\\\hline
\pronsrthreecone&---&ごと (goto)\\\hline
\end{prons}
\rindex{MAI}{毎}
\rindex{goto}{毎}

\begin{remprons}
\hooglig{Every} \ucomkun{goto} \comon{migration}.\\\hline
\end{remprons}

%{\middel}
% &  ()\\\hline
\begin{somme}
毎日&every + sun = \textit{daily}&マイ{\middel}ニチ \mbox{(MAI{\middel}NICHI)}\\\hline
毎週&every + week = \textit{weekly}&マイ{\middel}シュウ \mbox{(MAI{\middel}SH\={U})}\\\hline
毎月&every + moon (month) = \textit{montly}&マイ{\middel}つき \mbox{(MAI{\middel}tsuki)}\\\hline
毎年&every + year = \textit{yearly}&マイ{\middel}とし \mbox{(MAI{\middel}toshi)}\\\hline
毎朝&every + morning = \textit{every morning}&マイ{\middel}あさ \mbox{(MAI{\middel}asa)}\\\hline
毎回&every + rotate = \textit{each round}&マイ{\middel}カイ \mbox{(MAI{\middel}KAI)}\\\hline
毎時&every + time = \textit{hourly}&マイ{\middel}ジ \mbox{(MAI{\middel}JI)}\\\hline
毎秒&every + second = \textit{every second}&マイ{\middel}ビョウ \mbox{(MAI{\middel}BY\={O})}\\\hline
毎夕&every + evening = \textit{every evening}&マイ{\middel}ユウ \mbox{(MAI{\middel}Y\={U})}\\\hline
\notyet{毎晩}&every + late = \textit{every night}&マイ{\middel}バン \mbox{(MAI{\middel}BAN)}\\\hline
\notyet{毎度}&every + degree = \textit{each time}&マイ{\middel}ド \mbox{(MAI{\middel}DO)}\\\hline
\end{somme}

\begin{decomp}{5}
毎日&何時&に&起きます&か。\\\hline
マイニチ&なんジ&に&おきます&か\\\hline
every day&what time&に&get up&か\\\hline
\multicolumn{5}{|r|}{\textit{What time do you get up every day?}}\\\hline
\end{decomp}
\end{multicols}

\pagebreak
%%--------------------------------------------------------------------
%9--------------------------------------------------------------------
\begin{multicols}{2}
\begin{glyph}{Empty (空)}{empty}\label{glyph:empty}
Empty.  Sky.\\\hline
As in English, \textit{empty} can incorporate suggestions of \textit{uselessness} and \textit{futility}.
\end{glyph}
\gindex{Empty}{空}\strindex{8}{空}

\begin{comps}
\comprtwocone & 穴 + 工 \\\hline
\comprthreecone & Cave + Craft/Work \\\hline
\end{comps}

\begin{remglyph}
\hooglig{Empty} \remcom{caves} are \remcom{crafted} out.\\\hline
\end{remglyph}

\begin{prons}
\pronsrtwocone & \textbf{クウ (K\={U})} & \textbf{そら、から (sora, kara)} \\\hline
\rye{3}{クウ is by far the most common reading in compounds.\\\hline
そら means \textit{sky}.\\\hline
から occasionally appears in the first position with a meaning of \textit{empty}.}
\pronsrthreecone&---&あ (a)\\\hline
\rye{3}{The あ reading is used to make the verbs 空く and 空ける.  They can be thought of as interchangeable with the more common 開く and 開ける from Entry \ref{glyph:open} (p. \pageref{glyph:open}).}
\end{prons}
\rindex{KUU@K\={U}}{空}
\rindex{sora}{空}
\rindex{kara}{空}
\rindex{a}{空}

\begin{remprons}
\hooglig{Empty} \comkun{saw lumps} were \ucomkun{updated} to be \comon{cool}
and \comkun{caramelized}.\\\hline
\end{remprons}

%{\middel}
% &  ()\\\hline
\begin{somme}
\splitter{空く}{intr}&to be open&あ{\middel}く \mbox{(a{\middel}ku)}\\\hline
\splitter{空ける}{tr}&to open something&あ{\middel}ける \mbox{(a{\middel}keru)}\\\hline
空&\textit{sky}&そら \mbox{(sora)}\\\hline
青空&blue + empty = \textit{blue sky}&あお{\middel}ぞら \mbox{(ao{\middel}zora)}\\\hline
空手&empty + hand = \textit{karate}&から{\middel}て \mbox{(kara{\middel}te)}\\\hline
空間&empty + interval = \textit{space}&クウ{\middel}カン \mbox{(K\={U}{\middel}KAN)}\\\hline
空気&empty + spirit = \textit{air}&クウ{\middel}キ \mbox{(K\={U}{\middel}KI)}\\\hline
空中&empty + middle = \textit{mid-air}&クウ{\middel}チュウ \mbox{(K\={U}{\middel}CH\={U})}\\\hline
大空&large + open = \textit{firmament}&おお{\middel}ぞら \mbox{(\={o}{\middel}zora)}\\\hline
上空&upper + empty = \textit{the skies}&ジョウ{\middel}クウ \mbox{(J\={O}{\middel}K\={U})}\\\hline
\notyet{空港}&empty + port = \textit{Airport}&クウ{\middel}コウ \mbox{(K\={U}{\middel}K\={O})}\\\hline
\notyet{真空}&true + empty = \textit{vacuum; hollow}&シン{\middel}クウ \mbox{(SHIN{\middel}K\={U})}\\\hline
{\iscov}空き缶&empty + can = \textit{empty can}&あき{\middel}カン \mbox{(aki{\middel}KAN)}\\\hline
空き瓶{\notcov}&empty + bin = \textit{empty bin}&あき{\middel}ビン \mbox{(aki{\middel}BIN)}\\\hline
\end{somme}

\begin{decomp}{7}
朝&の&空&は&きれい&です&よ。\\\hline
あさ&の&そら&は&きれい&です&よ\\\hline
morning&の&skry&は&pretty&be/is&よ\\\hline
\multicolumn{7}{|r|}{\textit{The morning sky is lovely.}}\\\hline
\end{decomp}
\end{multicols}
%%--------------------------------------------------------------------
%10-------------------------------------------------------------------
\begin{multicols}{2}
\begin{glyph}{Town (町)}{town}\label{glyph:town}
Town.  One size smaller than a ``City'' (p. \pageref{glyph:city}).\\\hline
On occasion it can also refer to a \textit{street} or an \textit{alley}.
\end{glyph}
\gindex{Town}{町}\strindex{7}{町}

\begin{comps}
\comprtwocone&田 + 丁\\\hline
\comprthreecone&Rice field + Block (p. \pageref{glyph:block})\\\hline
\end{comps}

\begin{remglyph}
A \hooglig{town} of \remcom{rice-field} \remcom{blocks}.\\\hline
\end{remglyph}

\begin{prons}
\pronsrtwocone&\textbf{チョウ (CH\={O})}&\textbf{まち (machi)}\\\hline
\pronsrthreecone&---&---\\\hline
\end{prons}
\rindex{CHOU@CH\={O}}{町}
\rindex{machi}{町}

\begin{remprons}
\hooglig{Town} \comkun{marching} is such a \comon{chore}.\\\hline
\end{remprons}

%{\middel}
% &  ()\\\hline
\begin{somme}
町&\textit{town}&まち \mbox{(machi)}\\\hline
小町&small + town = \textit{belle; town-beauty}&こ{\middel}まち \mbox{(ko{\middel}machi)}\\\hline
町内&town + inside = \textit{in the town}&チョウ{\middel}ナイ \mbox{(CH\={O}{\middel}NAI)}\\\hline
町立&town + stand = \textit{established by the town}&チョウ{\middel}リツ \mbox{(CH\={O}{\middel}RITSU)}\\\hline
町会&town + meet = \textit{town assembly}&チョウ{\middel}カイ \mbox{(CH\={O}{\middel}KAI)}\\\hline
町中&town + middle = \textit{downtown}&まち{\middel}なか \mbox{(machi{\middel}naka)}\\\hline
町民&town + people = \textit{townspeople}&チョウ{\middel}ミン \mbox{CH\={O}{\middel}MIN}\\\hline
\notyet{町村}&town + village = \textit{towns and villages}&チョウ{\middel}ソン \mbox{(CH\={O}{\middel}SON)}\\\hline
\end{somme}

\begin{decomp}{7}
此の&町&の&ラーメン&は&有名&です。\\\hline
この&まち&の&ラーメン&は&ユウメイ&です\\\hline
this&town&の&ramen&は&famous&be/is\\\hline
\multicolumn{7}{|r|}{\textit{This town's ramen is famous.}}\\\hline
\end{decomp}
\end{multicols}
\pagebreak
%%--------------------------------------------------------------------
%11-------------------------------------------------------------------
\begin{multicols}{2}
\begin{glyph}{Dog (犬)}{dog}\label{glyph:dog}
Dog.
\end{glyph}
\gindex{Dog}{犬}\strindex{4}{犬}

\begin{comps}
\comprtwocone&犬\\\hline
\comprthreecone&Dog\\\hline
\end{comps}

\begin{remglyph}
Part of the 漢字 radicals (\#{94}).\\\hline
\end{remglyph}

\begin{prons}
\pronsrtwocone&\textbf{ケン (KEN)}&\textbf{いぬ (inu)}\\\hline
\pronsrthreecone&---&---\\\hline
\end{prons}
\rindex{KEN}{犬}
\rindex{inu}{犬}

\begin{remprons}
\hooglig{Dogs}, like \comkun{inuyasha}, have \comon{kennels}.\\\hline
{\warn}\textit{Inuyasha} is a famous アニメ character who is half-demon and half-human.\\\hline
\end{remprons}

%{\middel}
% &  ()\\\hline
\begin{somme}
犬&\textit{dog}&いぬ \mbox{(inu)}\\\hline
子犬&child + dog = \textit{puppy}&こ{\middel}いぬ \mbox{(ko{\middel}inu)}\\\hline
秋田犬&autumn + rice field + dog = \textit{an Akita dog}&あき{\middel}た{\middel}ケン \mbox{(aki{\middel}ta{\middel}KEN)}\\\hline
\notyet{愛犬}&love + dog = \textit{pet dog}&アイ{\middel}ケン \mbox{(AI{\middel}KEN)}\\\hline
\notyet{愛犬家}&love + dog + house = \textit{dog lover}&アイ{\middel}ケン{\middel}カ \mbox{(AI{\middel}KEN{\middel}KA)}\\\hline
\notyet{番犬}&order + dog = \textit{watch-dog}&バン{\middel}ケン \mbox{(BAN{\middel}KEN)}\\\hline
\notyet{野犬}&field + dog = \textit{stray dog}&ヤ{\middel}ケン \mbox{(YA{\middel}KEN)}\\\hline
\notyet{犬小屋}&dog + hut = \textit{kennel}&いぬ{\middel}ご{\middel}や \mbox{(inu{\middel}go{\middel}ya)}\\\hline
\notyet{軍犬}&army + dog = \textit{wardog}&グン{\middel}ケン \mbox{(GUN{\middel}KEN)}\\\hline
\end{somme}

\begin{decomp}{5}
古川さん&が&子犬&を&買いました。\\\hline
ふるかわさん&が&こいぬ&を&かいました\\\hline
Furukawa-san&が&puppy&を&bought\\\hline
\multicolumn{5}{|r|}{\textit{Furukawa-san bought a puppy.}}\\\hline
\end{decomp}
\end{multicols}
%%--------------------------------------------------------------------
%12-------------------------------------------------------------------
\begin{multicols}{2}
\begin{glyph}{Death (死)}{death}\label{glyph:death}
Death.
\end{glyph}
\gindex{Death}{死}\strindex{6}{死}

\begin{comps}
\comprtwocone&歹 + ヒ\\\hline
\comprthreecone&Death + Spoon/Ladle\\\hline
\end{comps}

\begin{remglyph}
\hooglig{Death} \remcom{spoon}.\\\hline
\hooglig{Death} on \remcom{one} \remcom{evening} from \remcom{spooning}.\\\hline
\end{remglyph}

\begin{prons}
\pronsrtwocone&\textbf{シ (SHI)}&\textbf{し (shi)}\\\hline
\pronsrthreecone&---&---\\\hline
\end{prons}
\rindex{SHI}{死}
\rindex{shi}{死}

\begin{remprons}
\hooglig{Death}\ldots  There is no \comon{shield} against or \comkun{shifting} away from it.\\\hline
\end{remprons}

%{\middel}
% &  ()\\\hline
\begin{somme}
\splitter{死ぬ}{intr}&\textit{to die}&し{\middel}ぬ \mbox{(shi{\middel}nu)}\\\hline
水死&water + death = \textit{drowning}&スイ{\middel}シ \mbox{(SUI{\middel}SHI)}\\\hline
死体&death + body = \textit{dead body}&シ{\middel}タイ \mbox{(SHI{\middel}TAI)}\\\hline
死火山&death + fire + mountain = \textit{extinct volcano}&シ{\middel}カ{\middel}ザン \mbox{(SHI{\middel}KA{\middel}ZAN)}\\\hline
死力&death + strength = \textit{desperate effort}&シリ{\middel}ョク \mbox{(SHI{\middel}RYOKU)}\\\hline
死者&death + individual = \textit{dead person; casualties}&シ{\middel}シャ \mbox{(SHI{\middel}SHA)}\\\hline
死語&death + word = \textit{extinct language}&シ{\middel}ゴ \mbox{(SHI{\middel}GO)}\\\hline
\notyet{必死}&certain + death = \textit{frantic; inevitable death}&ヒッ{\middel}シ \mbox{(HIS{\middel}SHI)}\\\hline
\notyet{死亡}&death + perish = \textit{death; mortality}&シ{\middel}ボウ \mbox{(SHI{\middel}B\={O})}\\\hline
\notyet{死産}&death + produce = \textit{stillbirth}&シ{\middel}ザン \mbox{(SHI{\middel}ZAN)}\\\hline
\notyet{死去}&death + depart = \textit{passing away}&シ{\middel}キョ \mbox{(SHI{\middel}KYO)}\\\hline
\notyet{死神}&death + god = \textit{god of death}&しに{\middel}がみ \mbox{(shini{\middel}gami)}\\\hline
\end{somme}

\begin{decomp}{4}
人間&は&必ず&死ぬ。\\\hline
ニンゲン&は&かならず&しぬ\\\hline
human&は&certainly&to die\\\hline
\multicolumn{4}{|r|}{\textit{Man is bound to die.}}\\\hline
\end{decomp}
\end{multicols}

\pagebreak
%%--------------------------------------------------------------------
%13-------------------------------------------------------------------
\begin{multicols}{2}
\begin{glyph}{Go (行)}{go}\label{glyph:go}
\textit{Go} is the primary meaning of this character.\\\hline
Related ideas are \textit{carrying out} or \textit{doing}.\\\hline
Far less common is the secondary sense of a \textit{stroke} or \textit{line} of text.
\end{glyph}
\gindex{Go}{行}\strindex{6}{行}

\begin{comps}
\comprtwocone&彳 + 一 + 丁\\\hline
\comprthreecone&Step + One + Block (p. \pageref{glyph:block})\\\hline
\end{comps}

\begin{remglyph}
\hooglig{Go} \remcom{step} on \remcom{one} \remcom{block}.\\\hline
\end{remglyph}

\begin{prons}
\pronsrtwocone&\textbf{コウ (K\={O})}&\textbf{い、ゆ (i, yu)}\\\hline
\rye{3}{For the general sense of the verb \textit{to go}, いく is used.\\\hline
ゆき is used in a few specific compounds, or as a suffix meaning \textit{bound for}.}
\pronsrthreecone&ギョウ、アン (GY\={O}, AN)&おこな (okona)\\\hline
\rye{3}{ギョウ is found primarily in compounds relating to this 漢字's secondary meaning of \textit{line} or \textit{stroke}.}
\end{prons}
\rindex{KOU@K\={O}}{行}
\rindex{i}{行}
\rindex{yu}{行}
\rindex{GYOU@GY\={O}}{行}
\rindex{AN}{行}
\rindex{okona}{行}

\begin{remprons}
\hooglig{Go} \comon{cause} a \ucomon{long-gyone} \comkun{interesting} \comkun{yule} in \ucomkun{a corner} with your \ucomon{uncle}.\\\hline
{\warn}\textit{Yule} is a festival historically observed by the Germanic peoples.\\\hline
\end{remprons}

%{\middel}
% &  ()\\\hline
\begin{somme}
\rye{3}{行 can also act like 買, but on rare occasions it incorporates ゆく as opposed to ゆき.}
行く&\textit{to go}&い{\middel}く \mbox{(i{\middel}ku)}\\\hline
行き先&go + precede = \textit{destination}&ゆ{\middel}き{\wsplit}さき \mbox{(yu{\middel}ki{\wsplit}saki)}\\\hline
東京行&east + capital + go = \textit{bound for Tokyo}&トウ{\middel}キョウ{\wsplit}ゆ{\middel}き \mbox{(T\={O}{\middel}KY\={O}{\wsplit}yu{\middel}ki)}\\\hline
夜行&night + go = \textit{night travel}&ヤ{\middel}コウ \mbox{(YA{\middel}K\={O})}\\\hline
歩行&walk + go = \textit{walking}&ホ{\middel}コウ \mbox{(HO{\middel}K\={O})}\\\hline
歩行者&walk + go + individual = \textit{pedestrian}&ホ{\middel}コウ{\middel}シャ \mbox{(HO{\middel}K\={O}{\middel}SHA)}  \\\hline
\notyet{行方不明}&whereabouts + unclear = \textit{missing}&ゆく{\middel}え{\wsplit}フ{\middel}メイ \mbox{(yuku{\middel}e{\wsplit}FU{\middel}MEI)}\\\hline
\notyet{急行}&hasten + go = \textit{hastening}&キュウ{\middel}コウ \mbox{(KY\={U}{\middel}K\={O})}\\\hline
\end{somme}

\begin{decomp}{5}
今日&は&歩行者&が&多い。\\\hline
きょう&は&ホコウシャ&が&おおい\\\hline
today&は&pedestrian&が&many\\\hline
\multicolumn{5}{|r|}{\textit{There are many pedestrians today.}}\\\hline
\end{decomp}
\end{multicols}
%%--------------------------------------------------------------------
%14-------------------------------------------------------------------
\begin{multicols}{2}
\begin{glyph}{Take (取)}{take}\label{glyph:take}
Take.
\end{glyph}
\gindex{Take}{取}\strindex{8}{取}

\begin{comps}
\comprtwocone&耳 + 又\\\hline
\comprthreecone&Ear + Again\\\hline
\end{comps}

\begin{remglyph}
\hooglig{Take} by the \remcom{ear} \remcom{again}.\\\hline
\end{remglyph}

\begin{prons}
\pronsrtwocone&---&\textbf{と (to)}\\\hline
\pronsrthreecone&シュ (SHU)&---\\\hline
\end{prons}
\rindex{to}{取}
\rindex{SHU}{取}

\begin{remprons}
\hooglig{Take} on the \comkun{torrential} \ucomon{shooting}.\\\hline
\end{remprons}

%{\middel}
% &  ()\\\hline
\begin{somme}
取る&\textit{to take}&と{\middel}る \mbox{(to{\middel}ru)}\\\hline
取引&take + pull = \textit{transactions}&とり{\wsplit}ひ{\middel}き \mbox{(tori{\wsplit}hi{\middel}ki)}\\\hline
\splitter{取り上げる}{tr}&take + upper = \textit{to pick up}&と{\middel}り{\wsplit}あ{\middel}げる \mbox{(to{\middel}ri{\wsplit}a{\middel}geru)}\\\hline
\splitter{取り入れる}{tr}&take + enter = \textit{to take in}&と{\middel}り{\wsplit}い{\middel}れる \mbox{(to{\middel}ri{\wsplit}i{\middel}reru)}\\\hline
\splitter{取り出す}{tr}&take + exit = \textit{to take out}&と{\middel}り{\wsplit}だ{\middel}す \mbox{(tori{\wsplit}dasu)}\\\hline
\splitter{気取る}{intr}&spirit + take = \textit{to put on airs}&キ{\wsplit}ど{\middel}る \mbox{(KI{\wsplit}do{\middel}ru)}\\\hline
\splitterny{受取る}{tr}&receive + take = \textit{to receive}&うけ{\wsplit}と{\middel}る \mbox{(uke{\wsplit}to{\middel}ru)}\\\hline
\notyet{受取り}&receive + take = \textit{receiving; receipt}&うけ{\wsplit}と{\middel}り \mbox{(uke{\wsplit}to{\middel}ri)}\\\hline
\end{somme}
\end{multicols}

\begin{decomp}{6}
良し、&此れ&で&引取&は&まとまった。\\\hline
よし&これ&で&とりひき&は&まとまった\\\hline
fine&this&で&transactions&は&to be settled\\\hline
\multicolumn{6}{|r|}{\textit{All right.  It's a deal.}}\\\hline
\end{decomp}

\pagebreak
%%--------------------------------------------------------------------
%15-------------------------------------------------------------------
\begin{multicols}{2}
\begin{glyph}{Village (村)}{village}\label{glyph:village}
Village.\\\hline
This character implies something smaller than a city or town.
\end{glyph}
\gindex{Village}{村}\strindex{7}{村}

\begin{comps}
\comprtwocone&木 + 寸\\\hline
\comprthreecone&Tree + Tiny\\\hline
\end{comps}

\begin{remglyph}
\hooglig{Village} \remcom{trees} are \remcom{tiny}.\\\hline
\end{remglyph}

\begin{prons}
\pronsrtwocone&\textbf{ソン (SON)}&\textbf{むら (mura)}\\\hline
\pronsrthreecone&---&---\\\hline
\end{prons}
\rindex{SON}{村}
\rindex{mura}{村}

\begin{remprons}
\hooglig{Village} \comon{songs} depicted in the \comkun{mural}.\\\hline
\end{remprons}

%{\middel}
% &  ()\\\hline
\begin{somme}
村&\textit{village}&むら \mbox{(mura)}\\\hline
村々&village $(\times 2)$ = \textit{villages}&むら{\middel}むら \mbox{(mura{\middel}mura)}\\\hline
村人&village + person = \textit{villager}&むら{\middel}びと \mbox{(mura{\middel}bito)}\\\hline
山村&mountain + village = \textit{mountain village}&サン{\middel}ソン \mbox{(SAN{\middel}SON)}\\\hline
村会&village + meet = \textit{village assembly}&ソン{\middel}カイ \mbox{(SON{\middel}KAI)}\\\hline
入村&enter + village = \textit{moving to a village}&ニュウ{\middel}ソン \mbox{(NY\={U}{\middel}SON)}\\\hline
\notyet{村民}&village + people = \textit{villager}&ソン{\middel}ミン \mbox{(SON{\middel}MIN)}\\\hline
\notyet{無医村}&nil + doctor + village = \textit{doctorless village}&ム{\middel}イ{\middel}ソン \mbox{(MU{\middel}I{\middel}SON)}\\\hline
\notyet{選手村}&athlete + village = \textit{Olympic village}&セン{\middel}シュ{\middel}むら \mbox{(SEN{\middel}SHU{\middel}mura)}\\\hline
\end{somme}

\begin{decomp}{5}
村人&は&小道&に&立っています。\\\hline
むらびと&は&こみち&に&たっています\\\hline
villager&は&path&に&is standing\\\hline
\multicolumn{5}{|r|}{\textit{A villager is standing in the path.}}\\\hline
\end{decomp}
\end{multicols}
%%--------------------------------------------------------------------
%16-------------------------------------------------------------------
\begin{multicols}{2}
\begin{glyph}{Character (字)}{character}\label{glyph:character}
Character in the sense of a \textit{Chinese character} or \textit{letter of an alphabet}.
\end{glyph}
\gindex{Character}{字}\strindex{6}{字}

\begin{comps}
\comprtwocone&宀 + 子\\\hline
\comprthreecone&Roof + Child\\\hline
\end{comps}

\begin{remglyph}
\hooglig{Character} under a \remcom{roof} is \remcom{childish}.\\\hline
\end{remglyph}

\begin{prons}
\pronsrtwocone&\textbf{ジ (JI)}&---\\\hline
\pronsrthreecone&---&あざ (aza)\\\hline
\end{prons}
\rindex{JI}{字}
\rindex{aza}{字}

\begin{remprons}
\hooglig{Characters} in the \ucomkun{azalea} \comon{jig}.\\\hline
{\warn}\textit{Azaleas} are flowering shrubs in the genus \textit{Rhododendron}.\\\hline
\end{remprons}

%{\middel}
% &  ()\\\hline
\begin{somme}
十字&ten + character = \textit{a cross}&ジュウ{\middel}ジ \mbox{(J\={U}{\middel}JI)}\\\hline
赤十字&red + ten + character = \textit{The Red Cross}&セキ{\middel}ジュウ{\middel}ジ \mbox{(SEKI{\middel}J\={U}{\middel}JI)}\\\hline
字引&character + pull = \textit{dictionary}&ジ{\middel}びき \mbox{(JI{\middel}biki)}\\\hline
文字&literature + character = \textit{letter; character}&モ{\middel}ジ \mbox{(MO{\middel}JI)}\\\hline
赤字&red + character = \textit{being in the red}&あか{\middel}ジ \mbox{(aka{\middel}JI)}\\\hline
大文字&large + character = \textit{capital letter}&おお{\middel}モ{\middel}ジ \mbox{(\={o}{\middel}MO{\middel}JI)}\\\hline
\notyet{黒字}&black + character = \textit{being in the black}&くろ{\middel}ジ \mbox{(kuro{\middel}JI)}\\\hline
\notyet{漢字}&China + character = \textit{kanji}&カン{\middel}ジ \mbox{(KAN{\middel}JI)}\\\hline
\notyet{数字}&number + character = \textit{number; digit}&スウ{\middel}ジ \mbox{(S\={U}{\middel}JI)}\\\hline
\notyet{文字通り}&character + passage = \textit{literally; to the letter}&モ{\middel}ジ{\wsplit}とお{\middel}り \mbox{(MO{\middel}JI{\wsplit}t\={o}{\middel}ri)}\\\hline
\end{somme}
\end{multicols}

\begin{decomp}{9}
多く&の&国&で&は&赤十字&は&大切&です。\\\hline
おおく&の&くに&で&は&セイジュウジ&は&タイセツ&です\\\hline
many&の&countries&で&は&Red Cross&は&important&be/is\\\hline
\multicolumn{9}{|r|}{\textit{The Red Cross is important in many countries.}}\\\hline
\end{decomp}

\pagebreak
%%--------------------------------------------------------------------
%17-------------------------------------------------------------------
\begin{multicols}{2}
\begin{glyph}{Tool (具)}{tool}\label{glyph:tool}
Tool.  Equipment.
\end{glyph}
\gindex{Tool}{具}\strindex{8}{具}

\begin{comps}
\comprtwocone&目 + 丌\\\hline
\comprthreecone&Eye + Table\\\hline
\end{comps}

\begin{remglyph}
\hooglig{Tools} for your \remcom{eyes} are on the \remcom{table}.\\\hline
\end{remglyph}

\begin{prons}
\pronsrtwocone&\textbf{グ (GU)}&---\\\hline
\pronsrthreecone&---&---\\\hline
\end{prons}
\rindex{GU}{具}

\begin{remprons}
\hooglig{Tool} \comon{goofing}.\\\hline
\end{remprons}

%{\middel}
% &  ()\\\hline
\begin{somme}
道具&road + tool = \textit{tool}&ドウ{\middel}グ \mbox{(D\={O}{\middel}GU)}\\\hline
家具&house + tool = \textit{furniture}&カ{\middel}グ \mbox{(KA{\middel}GU)}\\\hline
具体&tool + body = \textit{definite/concrete}&グ{\middel}タイ \mbox{(GU{\middel}TAI)}\\\hline
具体化&concrete + change = \textit{embodiment}&グ{\middel}タイ{\middel}カ \mbox{(GU{\middel}TAI{\middel}KA)}\\\hline
馬具&horse + tool = \textit{harness}&バ{\middel}グ \mbox{(BA{\middel}GU)}\\\hline
雨具&rain + tool = \textit{rain gear}&あま{\middel}グ \mbox{(ama{\middel}GU)}\\\hline
夜具&night + tool = \textit{bedding; bed clothes}&ヤ{\middel}グ \mbox{(YA{\middel}GU)}\\\hline
工具&craft + tool = \textit{tool; implement}&コウ{\middel}グ \mbox{(K\={O}{\middel}GU)}\\\hline
\notyet{具合}&tool + match = \textit{condition}&グ{\wsplit}あ{\middel}い \mbox{(GU{\wsplit}a{\middel}i)}\\\hline
\notyet{絵の具}&picture + tool = \textit{paint}&エ{\middel}の{\middel}グ \mbox{(E{\middel}no{\middel}GU)}\\\hline
\end{somme}

\begin{decomp}{7}
その&家具&は&母&の&物&です。\\\hline
その&カグ&は&はは&の&もの&です\\\hline
that&furniture&は&mother&の&thing&be/is\\\hline
\multicolumn{7}{|r|}{\textit{The furniture belongs to my mother.}}\\\hline
\end{decomp}
\end{multicols}
%%--------------------------------------------------------------------
%18-------------------------------------------------------------------
\begin{multicols}{2}
\begin{glyph}{Pour (注)}{pour}\label{glyph:pour}
Pour.\\\hline
\textit{Directing} something toward an object (a gaze or one's attention, for instance).
\end{glyph}
\gindex{Pour}{注}\strindex{8}{注}

\begin{comps}
\comprtwocone&氵 + 主\\\hline
\comprthreecone&Water + Primary (p. \pageref{glyph:primary})\\\hline
\end{comps}

\begin{remglyph}
\hooglig{Pour} \remcom{water} \remcom{primarily}.\\\hline
\end{remglyph}

\begin{prons}
\pronsrtwocone&\textbf{チュウ (CH\={U})}&\textbf{そそ (soso)}\\\hline
\pronsrthreecone&---&---\\\hline
\end{prons}
\rindex{CHUU@CH\={U}}{注}
\rindex{soso}{注}

\begin{remprons}
\hooglig{Pour} \comkun{so-so} while \comon{chewing}.\\\hline
\end{remprons}

%{\middel}
% &  ()\\\hline
\begin{somme}
\splitter{注ぐ}{tr}&\textit{to pour}&そそ{\middel}ぐ \mbox{(soso{\middel}gu)}\\\hline
注意&pour + mind = \textit{attention}&チュウ{\middel}イ \mbox{(CH\={U}{\middel}I)}\\\hline
不注意&not + attention = \textit{carelessness}&フ{\middel}チュウ{\middel}イ \mbox{(FU{\middel}CH\={U}{\middel}I)}\\\hline
注文書&order + write = \textit{written order}&チュウ{\middel}モン{\middel}ショ \mbox{(CH\={U}{\middel}MON{\middel}SHO)}\\\hline
注意書き&attention + write = \textit{notes; instructions}&チュウ{\middel}イ{\wsplit}が{\middel}き \mbox{(CH\={U}{\middel}I{\wsplit}ga{\middel}ki)}\\\hline
要注意&essential + attention = \textit{need for caution}&ヨウ{\middel}チュウ{\middel}イ \mbox{(Y\={O}{\middel}CH\={U}{\middel}I)}\\\hline
注水&pour + water = \textit{watering}&チュウ{\middel}スイ \mbox{(CH\={U}{\middel}SUI)}\\\hline
注入&pour + enter = \textit{infusion}&チュウ{\middel}ニュウ \mbox{(CH\={U}{\middel}NY\={U})}\\\hline
注目&pour + eye = \textit{notice}&チュウ{\middel}モク \mbox{(CH\={U}{\middel}MOKU)}\\\hline
\notyet{発注}&discharge + pour = \textit{ordering}&ハッ{\middel}チュウ \mbox{(HAC{\middel}CH\={U})}\\\hline
\notyet{注文品}&order + goods = \textit{ordered goods}&チュウ{\middel}モン{\middel}ヒン \mbox{(CH\={U}{\middel}MON{\middel}HIN)}\\\hline
\end{somme}

\begin{decomp}{5}
犬&に&注意&して&下さい。\\\hline
いぬ&に&チュウイ&して&ください\\\hline
dog&に&attention&do&please\\\hline
\multicolumn{5}{|r|}{\textit{Watch out for the dog.}}\\\hline
\end{decomp}
\end{multicols}

\pagebreak
%%--------------------------------------------------------------------
%19-------------------------------------------------------------------
\begin{multicols}{2}
\begin{glyph}{Goods (品)}{goods}\label{glyph:goods}
This character expresses the idea of \textit{goods} or \textit{articles} in the majority of compounds in which it appears.\\\hline
It can occasionally have the more abstract notion of \textit{refinement}.
\end{glyph}
\gindex{Goods}{品}\strindex{9}{品}

\begin{comps}
\comprtwocone&品\\\hline
\comprthreecone&Goods\\\hline
\end{comps}

\begin{remglyph}
An uncommon 漢字 radical.\\\hline
\end{remglyph}

\begin{prons}
\pronsrtwocone&\textbf{ヒン (HIN)}&\textbf{しな (shina)}\\\hline
\pronsrthreecone&---&---\\\hline
\end{prons}
\rindex{HIN}{品}
\rindex{shina}{品}

\begin{remprons}
The \hooglig{goods} \comkun{she nudged} at was a \comon{hint}.\\\hline
\end{remprons}

%{\middel}
% &  ()\\\hline
\begin{somme}
品&\textit{goods}&しな \mbox{(shina)}\\\hline
品物&goods + thing = \textit{merchandise}&しな{\middel}もの \mbox{(shina{\middel}mono)}\\\hline
品切れ&goods + cut = \textit{``out of stock''}&しな{\wsplit}ぎ{\middel}れ \mbox{(shina{\wsplit}gi{\middel}re)}\\\hline
上品&upper + goods = \textit{elegant}&ジョウ{\middel}ヒン \mbox{(J\={O}{\middel}HIN)}\\\hline
下品&lower + goods = \textit{vulgar}&ゲ{\middel}ヒン \mbox{(GE{\middel}HIN)}\\\hline
中古品&middle + old + goods = \textit{secondhand goods}&チュウ{\middel}コ{\middel}ヒン \mbox{(CH\={U}{\middel}KO{\middel}HIN)}\\\hline
手品&hand + goods = \textit{magic illusion}&て{\middel}じな \mbox{(te{\middel}jina)}\\\hline
\notyet{作品}&make + goods = \textit{a work (book, film, \ldots)}&サク{\middel}ヒン \mbox{(SAKU{\middel}HIN)}\\\hline
\notyet{食品}&eat + goods = \textit{food products}&ショク{\middel}ヒン \mbox{(SHOKU{\middel}HIN)}\\\hline
\notyet{品質}&goods + quality = \textit{quality}&ヒン{\middel}シツ \mbox{(HIN{\middel}SHITSU)}\\\hline
\notyet{部品}&section + goods = \textit{parts}&ブ{\middel}ヒン \mbox{(BU{\middel}HIN)}\\\hline
\end{somme}

\begin{decomp}{4}
此の&品物&は&安い。\\\hline
この&しなもの&は&やすい\\\hline
this&article&は&cheap\\\hline
\multicolumn{4}{|r|}{\textit{This article is cheap.}}\\\hline
\end{decomp}
\end{multicols}
%%--------------------------------------------------------------------
%20-------------------------------------------------------------------
\begin{multicols}{2}
\begin{glyph}{Hamlet (里)}{hamlet}\label{glyph:hamlet}
Hamlet.  Here we have the fourth and smallest of our populated units.\\\hline
This character can also signify the \textit{countryside}.
\end{glyph}
\gindex{Hamlet}{里}\strindex{7}{里}

\begin{comps}
\comprtwocone&里\\\hline
\comprthreecone&Village/3.93 km/2.44 miles\\\hline
\end{comps}

\begin{remglyph}
Part of the 漢字 radicals (\#{166}).\\\hline
\end{remglyph}

\begin{prons}
\pronsrtwocone&---&\textbf{さと (sato)}\\\hline
\pronsrthreecone&リ (RI)&---\\\hline
\end{prons}
\rindex{sato}{里}
\rindex{RI}{里}

\begin{remprons}
\hooglig{Hamlet} caused a \comkun{subtle} \ucomon{leap} in my heart.\\\hline
{\warn}\textit{Hamlet} is a tragedy written by Shakespeare.\\\hline
\end{remprons}

%{\middel}
% &  ()\\\hline
\begin{somme}
里&\textit{hamlet}&さと \mbox{(sato)}\\\hline
山里&mountain + hamlet = \textit{mountain hamlet}&やま{\middel}さと \mbox{(yama{\middel}sato)}\\\hline
古里&old + hamlet = \textit{the old hometown}&ふる{\middel}さと \mbox{(furu{\middel}sato)}\\\hline
里人&hamlet + person = \textit{country folk}&さと{\middel}びと \mbox{(sato{\middel}bito)}\\\hline
里心&hamlet + heart = \textit{homesickness}&さと{\middel}ごころ \mbox{(sato{\middel}gokoro)}\\\hline
里親&hamlet + parents = \textit{foster parents}&さと{\middel}おや \mbox{(sato{\middel}oya)}\\\hline
里子&hamlet + child = \textit{foster child}&さと{\middel}ご \mbox{(sato{\middel}go)}\\\hline
\notyet{人里離れた}&human habitation + separate = \textit{lonely/remote place}&ひと{\middel}ざと{\wsplit}はな{\middel}れた \mbox{(hito{\middel}zato{\wsplit}hana{\middel}reta)}\\\hline
\notyet{里帰り}&hamlet + homeward = \textit{returning home}&さと{\wsplit}かえ{\middel}り \mbox{(sato{\wsplit}kae{\middel}ri)}\\\hline
\end{somme}

\begin{decomp}{5}
先週&は&古里&へ&行きました。\\\hline
センシュウ&は&ふるさと&へ&いきました\\\hline
last week&は&old hometown&へ&went\\\hline
\multicolumn{5}{|r|}{\textit{Last week I went to my old hometown.}}\\\hline
\end{decomp}
\end{multicols}
%%--------------------------------------------------------------------
%%--------------------------------------------------------------------
\pagebreak
